% ============================================================
%  ANSWER KEY — Practice 5: Combinatorics
% ============================================================

\section*{Վարժություն 5. Կոմբինատորիկա. Հիմնական կանոններ}

% ------------------------------------------------------------
\subsection*{Բաժին 1. Գումարի և արտադրյալի կանոններ}

\subsubsection*{Խնդիր 1}
0-ից 1000 ամբողջ թվերի քանակը, որոնք պարունակում են գոնե մեկ 6 թվանշան:

Լրացում. հաշվել այն ամբողջ թվերը, որոնք \emph{չունեն} 6 թվանշան:
\begin{itemize}
  \item 1-նիշ թվեր (0--9): 9 հատ չունեն 6 (բոլորը, բացի 6-ից):
  \item 2-նիշ թվեր (10--99): առաջին թվանշանը $\{1,\dots,9\}\setminus\{6\}$-ից՝ 8 հատ, երկրորդ թվանշանը $\{0,\dots,9\}\setminus\{6\}$-ից՝ 9 հատ: Ընդամենը $8\cdot9=72$:
  \item 3-նիշ թվեր (100--999): $8\cdot9\cdot9=648$:
  \item 1000 թիվը չի պարունակում 6: $+1$ հատ:
\end{itemize}
Առանց վեցերի ընդհանուր քանակը՝ $9+72+648+1=730$:
0-ից 1000 բոլոր ամբողջ թվերը՝ $1001$:

$\boxed{1001 - 730 = 271}$

\subsubsection*{Խնդիր 2}
Նույն խնդիրն է, ինչ Խնդիր 1-ը (վարժությունների թերթիկում կրկնություն է):

$\boxed{271}$

\subsubsection*{Խնդիր 3}
Ընտրել երկու տարբեր առարկաների գրքեր 15 ինֆորմատիկայի, 12 մաթեմատիկայի, 10 ֆիզիկայի գրքերից:

Ըստ գումարի/արտադրյալի կանոնների՝
$15\cdot12 + 15\cdot10 + 12\cdot10 = 180+150+120 = \boxed{450}$:

\subsubsection*{Խնդիր 4}
Կենտ թվեր 100-ից 1000 միջակայքում (այսինքն՝ 3-նիշ կենտ թվեր 101, 103, \ldots, 999):

Վերջին թվանշանը՝ 5 ընտրություն ($\{1,3,5,7,9\}$):
Առաջին թվանշանը՝ 9 ընտրություն ($1$--$9$):
Միջին թվանշանը՝ 10 ընտրություն ($0$--$9$):

$\boxed{9 \times 10 \times 5 = 450}$

\subsubsection*{Խնդիր 5}
700-ից փոքր բնական թվեր ($\{1,\dots,699\}$).

(a) Բաժանվում են 5-ի. $\lfloor 699/5\rfloor = \boxed{139}$:

(b) Բաժանվում են 3-ի. $\lfloor 699/3\rfloor = \boxed{233}$:

(c) Բաժանվում են և՛ 3-ի, և՛ 5-ի (այսինքն՝ 15-ի). $\lfloor 699/15\rfloor = \boxed{46}$:

\subsubsection*{Խնդիր 6}
Երևան $\to$ Մոսկվա: 15 չվերթ: Մոսկվա $\to$ Լոնդոն: 20 չվերթ:
Ըստ արտադրյալի կանոնի՝ $\boxed{15 \times 20 = 300}$:

\subsubsection*{Խնդիր 7}
4-նիշ թվեր, որոնք ունեն առնվազն երկու միանման թվանշաններ:

Բոլոր 4-նիշ թվերը՝ $9\times 10^3 = 9000$:
4-նիշ թվեր, որոնց բոլոր թվանշանները տարբեր են. առաջին թվանշանը՝ 9 ընտրություն, երկրորդը՝ 9, երրորդը՝ 8, չորրորդը՝ 7, ստացվում է $9\times 9\times 8\times 7 = 4536$:

$\boxed{9000 - 4536 = 4464}$

\subsubsection*{Խնդիր 8}
7-նիշ թվեր, որոնց առաջին և վերջին թվանշանները տարբեր են:

Բոլոր 7-նիշ թվերը՝ $9\times 10^6$:
7-նիշ թվեր, որոնց առաջին = վերջին թվանշանը. առաջին թվանշանը՝ 9 ընտրություն, վերջին թվանշանը ֆիքսված է (հավասար է առաջինին՝ 1 ընտրություն), միջին 5 թվանշանները՝ $10^5$ յուրաքանչյուրը: Ընդամենը՝ $9\times 10^5$:

Պատասխան՝ $9\times 10^6 - 9\times 10^5 = 9\times 10^5(10-1)=9\times 10^5\times 9$:

$\boxed{9\times 10^6 - 9\times 10^5 = 8{,}100{,}000}$

\subsubsection*{Խնդիր 9}
6-նիշ թվեր, որտեղ հարևան թվանշանները տարբեր են (թվանշանները 0--9, առաջին թվանշանը $\neq 0$):

Դիրք 1: 9 ընտրություն (1--9): Հաջորդ յուրաքանչյուր դիրք՝ 9 ընտրություն (ցանկացած թվանշան, բացի նախորդից):

$\boxed{9 \times 9^5 = 9^6 = 531{,}441}$

\subsubsection*{Խնդիր 10}
7-նիշ թվեր. զույգ դիրքերում (2,4,6) կենտ թվանշաններ են, կենտ դիրքերում (1,3,5,7)՝ զույգ թվանշաններ:

Դիրքեր 2,4,6: 5 կենտ թվանշան յուրաքանչյուրում $\Rightarrow 5^3$:
Դիրք 1 (առաջին, պետք է լինի զույգ և $\neq 0$): 4 ընտրություն $\{2,4,6,8\}$:
Դիրքեր 3,5,7: 5 զույգ թվանշան յուրաքանչյուրում $\Rightarrow 5^3$:

$\boxed{4 \times 5^3 \times 5^3 = 4 \times 15{,}625 = 62{,}500}$

% ------------------------------------------------------------
\subsection*{Բաժին 2. Տեղափոխություններ և կարգավորություններ}

\subsubsection*{Խնդիր 1}
Արտապատկերումներ 6-տարրանոց բազմությունից 3-տարրանոց բազմություն. 6 տարրերից յուրաքանչյուրը արտապատկերվում է 3 պատկերներից մեկին:

$\boxed{3^6 = 729}$

\subsubsection*{Խնդիր 2}
Ինյեկտիվ արտապատկերումներ 6-տարրանոց բազմությունից 3-տարրանոց բազմություն:
Քանի որ $6 > 3$, ըստ Դիրիխլեի սկզբունքի (Pigeonhole Principle), ոչ մի ինյեկցիա գոյություն չունի:

$\boxed{0}$

\subsubsection*{Խնդիր 3}
5 տղա և 6 աղջիկ շարքում; տղաները զբաղեցնում են առաջին 5 դիրքերը:

Դասավորել 5 տղաներին առաջին 5 դիրքերում՝ $5!$:
Դասավորել 6 աղջիկներին մնացած 6 դիրքերում՝ $6!$:

$\boxed{5! \times 6! = 120 \times 720 = 86{,}400}$

\subsubsection*{Խնդիր 4}
6 տղա և 7 աղջիկ ($13$ մարդ) շարքում; առաջին դիրքում տղա է:

Ընտրել տղային 1-ին դիրքի համար՝ 6 եղանակ: Դասավորել մնացած 12 մարդկանց՝ $12!$:

$\boxed{6 \times 12! = 6 \times 479{,}001{,}600 = 2{,}874{,}009{,}600}$

\subsubsection*{Խնդիր 5}
6 տղա և 7 աղջիկ շարքում; առաջին և վերջին դիրքերում տղաներ են:

Առաջին դիրք՝ 6 ընտրություն: Վերջին դիրք՝ 5 մնացած տղաներից: Միջին 11 դիրքերը՝ $11!$:

$\boxed{6 \times 5 \times 11! = 30 \times 39{,}916{,}800 = 1{,}197{,}504{,}000}$

\subsubsection*{Խնդիր 6}
5 տղա և 6 աղջիկ շարքում; տղաները զբաղեցնում են 5 \emph{հաջորդական} դիրքեր:

Դիտարկենք 5 տղաների խումբը որպես մեկ միավոր: 6 աղջիկների հետ միասին ունենք 7 օբյեկտ շարքում դասավորելու համար՝ $7!$ եղանակ: Խմբի ներսում 5 տղաները կարող են տեղափոխվել՝ $5!$ եղանակով:

$\boxed{7! \times 5! = 5040 \times 120 = 604{,}800}$

\subsubsection*{Խնդիր 7}
Դասավորել 9 տարբեր մարդկանց շարքում:

$\boxed{9! = 362{,}880}$

\subsubsection*{Խնդիր 8}
8-նիշ թվեր՝ կազմված 23122313 թվի թվանշաններից:

Թվանշաններ. 1-ը հանդիպում է 2 անգամ, 2-ը՝ 3 անգամ, 3-ը՝ 3 անգամ: Ընդամենը 8 թվանշան:

Կրկնություններով տեղափոխություններ՝ $\dfrac{8!}{2!\,3!\,3!} = \dfrac{40320}{2\cdot6\cdot6} = \dfrac{40320}{72}$:

$\boxed{\frac{8!}{2!\,3!\,3!} = 560}$

\subsubsection*{Խնդիր 9}
6 երկարությամբ բառեր $\{a, b, g, d, e, z\}$ այբուբենից (ձայնավորներ՝ $a, e$, բաղաձայններ՝ $b, g, d, z$), որտեղ երկու ձայնավորներ հարևան չեն:

6 երկարությամբ բոլոր բառերը՝ $6^6 = 46656$:

Բառեր, որտեղ ՉԿԱՆ երկու հարևան ձայնավորներ. Դիցուք $V_n, C_n$ = ձայնավորով/բաղաձայնով վերջացող բառերի քանակը: $V_n = 2C_{n-1}$, $C_n = 4(V_{n-1}+C_{n-1})$:

$V_1=2, C_1=4, V_2=8, C_2=24, V_3=48, C_3=128, V_4=256, C_4=704$,
$V_5=1408, C_5=3840, V_6=7680, C_6=20992$: Ընդամենը՝ $f(6)=28672$:

$\boxed{46656 - 28672 = 17984}$

\subsubsection*{Խնդիր 10}
\textbf{Ապացուցել.} 6 հոգուց բաղկացած ցանկացած խմբում կա՛մ կան 3 փոխադարձ ծանոթներ, կա՛մ 3 փոխադարձ անծանոթներ ($R(3,3)=6$):

\textit{Ապացույցի ուրվագիծ.}
Մոդելավորենք 6 մարդկանց որպես $K_6$-ի գագաթներ՝ կողերը ներկված կարմիր (ծանոթ) կամ կապույտ (անծանոթ): Վերցնենք որևէ $A$ գագաթ; այն ունի 5 կող, ուստի ըստ Դիրիխլեի սկզբունքի՝ առնվազն 3-ը ունեն նույն գույնը, ասենք՝ կարմիր, դեպի $B,C,D$ գագաթները: Եթե $B,C,D$-ի միջև որևէ կող կարմիր է, այն $A$-ի հետ կազմում է կարմիր եռանկյուն; հակառակ դեպքում $B,C,D$-ն կազմում են կապույտ եռանկյուն: \qed

\subsubsection*{Խնդիր 11}
Քանի՞ թիվ պետք է ընտրել $\{1,2,\dots,210\}$ բազմությունից՝ երաշխավորելու համար, որ կա գոնե մեկ զույգ թիվ:

$\{1,\dots,210\}$-ում կան 105 կենտ թվեր: Վատագույն դեպքում մենք ընտրում ենք բոլոր 105 կենտ թվերը: Եվս մեկ ընտրություն երաշխավորում է զույգ թիվ:

$\boxed{106}$

\subsubsection*{Խնդիր 12}
\textbf{Ապացուցել.} $7, 77, 777, \ldots$ հաջորդականության առաջին 2014 անդամների մեջ կա մեկը, որը բաժանվում է 2013-ի:

\textit{Ապացույցի ուրվագիծ.}
Դիցուք $a_k = \underbrace{77\ldots7}_{k}$: Դիտարկենք 2014 մնացորդները $a_k \bmod 2013$: Եթե որևէ մնացորդ 0 է, ապացուցված է: Հակառակ դեպքում 2014 մնացորդները ընկնում են $\{1,\dots,2012\}$ բազմության մեջ (2012 դաս), ուստի ըստ Դիրիխլեի սկզբունքի՝ երկու անդամ $a_i, a_j$ ($i<j$) ունեն նույն մնացորդը: Այդ դեպքում $2013 \mid a_j - a_i = a_{j-i}\cdot 10^i$: Քանի որ $\gcd(10,2013)=1$, ստանում ենք $2013\mid a_{j-i}$, և $1\le j-i\le 2013$: \qed